\section{Learning Event Rules from Perception}\label{sec:ecinduction}

The preceding sections build the machinery; this one puts all of it on an external benchmark at once. The task is symbolic event-rule induction over video surveillance data, and the perceptual predicate the induced rules depend on is itself learned---through the logic credit alone, with no perceptual supervision anywhere in the loop.

\subsection{The Problem}

The Event Calculus is the standard logical formalism for narrative reasoning over time. A \emph{fluent} is a time-varying property of the world---for instance, that two tracked people are \texttt{meeting}. Two event predicates say when a fluent changes: \texttt{initiatedAt(F,T)} holds when something at time \texttt{T} makes fluent \texttt{F} begin, and \texttt{terminatedAt(F,T)} when something makes it stop. A domain-independent axiom of inertia closes the loop: a fluent \emph{holds at} a time if it was initiated earlier and not terminated since. Reasoning is therefore cheap and generic once the initiation and termination rules are known; the whole difficulty sits in obtaining them.

Those rules are classically hand-written by a domain expert. The line of systems that learns them instead---inductive logic programming over Event-Calculus theories---learns them from a symbolic input vocabulary that somebody else has already computed. On the CAVIAR video benchmark this means the learner is handed a relation such as \texttt{close(P1,P2,T)}, already reduced from raw tracker coordinates by a hand-chosen distance threshold, alongside per-person activity states. Rule induction thus starts one layer above perception, and the quality of the perceptual layer is a fixed input rather than something the learning explains.

\subsection{What Is Learned Here}

In the runs reported below the proximity predicate is not given. It is \texttt{close\_nn}, a small network over the raw pair coordinates, and it is never shown a distance label, a threshold, or any target of its own. Its only gradient arrives through the induced rules: a candidate clause containing \texttt{close\_nn} in its body is scored by how well the resulting Event-Calculus theory explains the frame-level annotations, and that score is what propagates back to the network's weights.

Mechanically this is exactly the trainable rule bodies of \Cref{subsec:trainable-bodies}, carrying a real benchmark rather than a demonstration. The neural atom sits inside a rule body; the symbolic literals beside it act as hard join conditions that gate which groundings fire without carrying gradient; credit reaches the network through the same verified circuit and the same zero-copy export described there. Structure search and weight training are not two pipelines glued together---the rule learner's own objective is the network's loss.

\subsection{Protocol}

Three guards define the protocol, and all three are pre-registered---declared before the run they govern, with no configurations executed on the real data outside those declared.

\textbf{Leak-free scene-family splitting.} The distributed CAVIAR corpus used by the published Event-Calculus learners is a dump of 26 video segments, and it is not clean: an independent frame-level audit of every \texttt{meeting} transition event in the dump against the original ground-truth XML classifies 21 of 25 events as real, 3 as duplicates (two source videos appear in both the dump's train and its test file) and 1 as a splice artifact; one video alone contributes 855 duplicated gold pair-frames. Cross-validating over segments therefore still permits same-scene transfer. The clean protocol folds over \emph{scene families} recovered from the source XML instead, so no scene appears on both sides of a split.

\textbf{Recall-aware holdout.} Scoring candidates by held-out accuracy is useless at this class balance---roughly ten positive transitions against tens of thousands of rows---because every candidate sits on the all-false base-rate plateau, and the ranking provably inverts: the best real detector scores below the empty theory. Candidate selection uses per-fold held-out F1 instead.

\textbf{Permutation-null fit gate.} A candidate must beat what label shuffling alone achieves on its own fold. Each fold derives its own threshold from 1000 label permutations (seed 7) as the 95th percentile of the pool-maximum mean per-fold F1, computed per target and on the training side only. A theory whose clauses do not clear that bar is not returned: the search abstains and reports an empty theory rather than committing one. Abstention is a first-class outcome of this protocol, not a failure mode of it.

\subsection{Results}

\textbf{Perception learned through logic credit.} On the windowed-fold protocol with a direct per-timestep target, the search finds the same two-clause theory on every fold in both the relational and the neural mode:

\begin{lstlisting}
holdsAt_meeting(PP,T) :- both_inactive(PP,T), close*(PP,T).
holdsAt_meeting(PP,T) :- both_active(PP,T),   close*(PP,T).
\end{lstlisting}

\noindent where \texttt{close*} is the precomputed geometric predicate in the relational mode and the learned \texttt{close\_nn} in the neural mode. Held-out F1 per fold is $0.9215 / 0.6804 / 0.7695$ with the precomputed predicate (mean $0.7905$) and $0.9215 / 0.7918 / 0.7695$ with the learned one (mean $0.8276$). The learned detector reproduces ground-truth geometry exactly on two folds and exceeds it on the third, where a soft learned boundary survives the train/test geometry shift better than a hard threshold. This is the paper's actual claim on this benchmark---a perceptual predicate trained through nothing but symbolic credit standing in for hand-set geometry without loss---and it rests on no published comparison at all.

\textbf{Protocol-matched cross-validation.} Under the full Event-Calculus target, the combined 26-segment corpus (32{,}360 co-visible pair-frames, 1{,}833 gold \texttt{meeting} frames) was cross-validated ten-fold by video segment with every gate re-derived inside each fold, and scored by micro-aggregation over pooled per-fold counts, as the published evaluations do. The result is precision $0.6580$, recall $0.8271$, F1 $0.7329$ (1516 true positives, 788 false positives, 317 false negatives). The same initiation clause---\texttt{both\_active \& close}---is selected on all ten folds, each time over that fold's own permutation null; the termination theory is empty on all ten, so the fluent persists by inertia to the end of each co-visible pair-run.

\textbf{Comparison, not ranking.} This protocol matches the published Event-Calculus learners on the axis that matters most for the metric: F1 is computed on \texttt{holdsAt} inferred through inertia, on the same distributed corpus, with the same micro-aggregation. The published figures for \texttt{meeting} on that corpus are $0.735$ for hand-crafted rules, $0.782$ and $0.792$ for OLED (the latter as reported by its authors, the former as re-run in the WOLED paper) and $0.887$ for WOLED-ASP; they are tabulated beside ours in \Cref{sec:eval}. Our $0.733$ sits just under that group's lowest figure. We draw no ordering from any of these differences, in either direction, because five protocol deltas separate the rows and they do not all point the same way. (i)~The published runs cross-validate over windowed interpretations with subsampled negatives; ours folds over whole video segments. (ii)~They learn full initiation \emph{and} termination programs; the run above has a two-literal initiation clause and an empty termination theory, which costs 788 of 2{,}304 predicted-positive frames as false positives on the folds with detected interior terminations. (iii)~Their rule language admits longer bodies---bottom clauses averaging some fifteen literals against our two-literal cap. (iv)~Their settings are reported as the best among several parameter settings tried; ours are pre-registered. (v)~Both compute F1 on \texttt{holdsAt} through inertia, which is why the comparison is drawn at all. The honest contrast is one of protocols, not of systems.

\textbf{A second iteration, labeled as such.} Adding two pair-level transition relations to the Event-Calculus vocabulary---one for a pair crossing the proximity threshold outward, one for strictly growing distance---and re-running the identical folds, seeds and gates raises the result to precision $0.7357$, recall $0.8260$, F1 $0.7782$ (1514/544/319), with a termination theory now learned on seven of the ten folds. That figure falls inside the $0.735$--$0.887$ span the published systems occupy, but it must never be read as the headline: its vocabulary was chosen \emph{after} the first run's test-side failure mode was known, on the same folds. Everything else about it is pre-registered; the vocabulary is not. Against published estimates selected as the best of several tried, the smaller fully pre-registered number is the stronger scientific claim, which is why $0.733$ leads here and $0.7782$ appears only with this label attached.

\textbf{A negative result.} Substituting the learned \texttt{close\_nn} detector for the precomputed proximity predicate inside the Event-Calculus initiation search does not survive cross-validation: precision $0.1250$, recall $0.0005$, F1 $0.0011$ (1/7/1832) over the same ten folds. Only one fold commits a clause---the same \texttt{both\_active \& close\_nn} shape the geometric predicate selects there, and probed directly against held-out ground-truth proximity that clause scores precision $1.0$ at recall $0.27$, so the network has learned the geometric relation. The failure is one of scale in the selection rule, not of perception: the initiation search scores coverage against transition \emph{events}, a handful per fold, and a probability-thresholded gate newly covers fewer of them than a deterministic predicate does, so nine of ten folds reject the candidate for insufficient new coverage. We report this as a result rather than omitting it; the direct protocol above is where the learned detector currently pays off.

\textbf{The clean protocol abstains.} Re-run on the deduplicated, scene-family-grouped corpus, the \texttt{meeting} fluent yields no theory at all---both the Event-Calculus route and the direct route abstain on every fold, F1 $0.0$. The reason is visible in the data and is not a defect of the search. One scene family carries 1{,}323 of the 1{,}812 \texttt{meeting}-positive frames and, under honest grouping, sits in exactly one fold; the clause that explains those frames holds \emph{only} there, covering zero of the 489 positives on all eight remaining \texttt{meeting}-bearing segments. Training with that family learns a clause that transfers nowhere; training without it cannot learn the clause at all. With eleven observable \texttt{meeting} initiations and five \texttt{moving} ones in the clean corpus, there is not enough evidence for a cross-validated theory, and the fit gate says so. That is the gate working as designed---a rule learner declining to manufacture a theory---and it is a more interesting property of the system than any F1 it reports elsewhere. The machinery is not broken under this protocol: on \texttt{moving} the direct route still selects the canonical published rule shape on nine of ten folds (micro F1 $0.4450$), and the \texttt{meeting} clause, applied directly to the held-out scene fold rather than transferred to it, scores frame F1 $0.890$.

\begin{table}[t]
  \centering
  \small
  \setlength{\tabcolsep}{4pt}
  \begin{tabularx}{\columnwidth}{Lccc}
    \toprule
    \textbf{Run} & \textbf{P} & \textbf{R} & \textbf{F1}\\
    \midrule
    \multicolumn{4}{l}{\emph{Windowed folds, direct target (mean of 3 folds)}}\\
    \quad precomputed \texttt{close} & --- & --- & 0.7905\\
    \quad learned \texttt{close\_nn} & --- & --- & \textbf{0.8276}\\
    \addlinespace
    \multicolumn{4}{l}{\emph{Distributed corpus, 10-fold CV, EC + inertia}}\\
    \quad pre-registered & 0.6580 & 0.8271 & \textbf{0.7329}\\
    \quad + termination vocab.\ (2nd iter.) & 0.7357 & 0.8260 & 0.7782\\
    \quad learned \texttt{close\_nn} & 0.1250 & 0.0005 & 0.0011\\
    \addlinespace
    \multicolumn{4}{l}{\emph{Clean protocol (scene-family folds)}}\\
    \quad \texttt{meeting}, both routes & --- & --- & abstains\\
    \quad \texttt{moving}, EC route & --- & --- & abstains\\
    \quad \texttt{moving}, direct route & 0.5334 & 0.3817 & 0.4450\\
    \bottomrule
  \end{tabularx}
  \caption{Runs reported in this section. Micro-averages over pooled per-fold counts, except the windowed rows, which are fold means. ``Abstains'' is an empty theory returned by the permutation-null fit gate, not a zero score. Published figures for the same benchmark are tabulated separately in \Cref{sec:eval}.}
  \label{tab:ecinduction}
\end{table}

\subsection{What This Establishes, and What It Does Not}

It establishes that a perceptual predicate can be learned end-to-end through symbolic credit alone and stand in for hand-set geometry without loss on a real benchmark (mean F1 $0.8276$ against $0.7905$); that Event-Calculus rules can be induced under fully pre-registered gates and land within a couple of points of the range published dedicated Event-Calculus learners report on the corpus they share; that the benchmark corpus those numbers rest on carries a measurable duplication defect, established by a re-runnable audit; and that the fit gate abstains rather than fabricates when a leak-free split leaves the evidence too thin.

It does not establish a ranking against OLED, WOLED-ASP or hand-crafted rules, and none is claimed: the protocol deltas above are real, unquantified, and cut both ways, and the clean protocol has too few events for a cross-validated comparison to mean anything in either direction. It does not show that termination programs are learned in the pre-registered run---they are not; the empty termination theory is the dominant source of false positives. It does not show that the learned detector is a drop-in replacement for the geometric predicate everywhere; under Event-Calculus cross-validation it is not. And it is one benchmark in one domain. What it does show is that the structure search and the perception underneath it can be one differentiable object rather than two systems in sequence.
